\honest{Could Anything Be Right?}{Eliezer Yudkowsky}{2008}

Years ago I was convinced that I knew nothing about morality. For all I knew, morality might require the extermination of the human species, and if it did, I saw no virtue in opposing it.

That was ``massively mistaken.'' As Davidson observed, I say, if you believe that ``beavers'' live in deserts, are pure white and weigh 300 pounds, you have no beliefs about beavers at all. ``You must get at least some of your beliefs right, before the remaining ones can be wrong about anything.'' \nb{The beaver example, in this wording, is from Richard Rorty's 2003 newspaper essay on Davidson. Davidson's own statement is stronger: a belief needs ``endless true beliefs about the subject matter.''} So my belief that I had no information about morality was not consistent; I did think killing was probably wrong.

My plan then was to build a generic superintelligence that would reason out what was right. But by the no-free-lunch theorems, if you know nothing, no procedure is more likely than another to compute morality; a program finds ``morality'' only if something in its start state points it there.

If you know nothing about morality, what distinguishes ``I don't know what is right'' from ``I don't know what is wakalixes''? A scientist's idea of evidence can change, as when a reader of Popper meets the proofs behind Bayesian probability, but only because the brain already contained what the new argument appeals to. In the same way you can say, ``I'll know it when I see it.''

You may distrust intuitions that evolution produced. But if you discard everything evolution gave you, you discard your whole brain, including the intuitions that tell you evolution is a poor source of morality. And a morality that no human could ever be moved to adopt would be no more morality than a stone tablet saying ``Thou shalt murder.''

So you might know a little about morality: ``that which your brain would recognize as right is what you are talking about.'' \nb{Starting from one's own judgments and revising them is the method of reflective equilibrium; Rawls called such judgments ``provisional fixed points.''} It follows that a ghost of perfect emptiness might not agree with you, since it might be ``asking a different question.'' \nb{The main objection to views on which our starting points fix what moral words mean is that people who start differently would then talk past each other instead of disagreeing. The post draws the consequence for a ghost, not for humans whose starting points differ.}

And if some things are built into the meaning of ``morality,'' ``then why not accept other intuitions''? Why not accept that, other things equal, joy is preferable to sorrow, not as a personal preference but as part of the question ``What is truly right''? \nb{The earlier argument showed that some content must be built in. This step, made by a question, does not show which intuitions belong to the question rather than to its answers. The resulting view resembles Frank Jackson's moral functionalism, on which moral properties are whatever fits the platitudes of a folk morality refined by reflection.} Then you know rather a lot about morality, though nothing certain.
